\chapter{Hamiltonian Monte Carlo}
\label{ch:hmc}

% PDF pages 99--111 (book pages 87--98)

The key idea of Hamiltonian Monte Carlo (HMC) is to extend the state space and exploit some
properties of Hamiltonian dynamics to avoid the local behavior of RWMH and Langevin Monte
Carlo. We write our target as
\[
f(q) \propto \exp\!\bigl(-V(q)\bigr), \qquad q \in \R^d,
\]
and introduce a distribution $f^H$ in extended space $\R^{2d}$
\begin{equation}
\label{eq:8.1}
\begin{aligned}
f^H(q,p) &:= \frac{1}{Z}\exp\!\bigl(-V(q)\bigr)\exp\!\bigl(-K(p)\bigr) \\
          &= \frac{1}{Z}\exp\!\bigl(-H(q,p)\bigr),
\end{aligned}
\end{equation}
where $K(p)$ satisfies $\int_{\R^d} \exp(-K(p))\,dp < \infty$ and
\[
H(q,p) := V(q) + K(p)
\]
is called the Hamiltonian. We call the measure $\mu^H$ with Lebesgue density $f^H$ the
canonical measure of the Hamiltonian $H$. All the ingredients just introduced have both a
physical and a statistical interpretation.
\begin{itemize}
  \item $q \in \R^d$ is interpreted as position $\equiv$ variables of interest.
  \item $p \in \R^d$ is interpreted as momentum $\equiv$ artificial variables that are not
    part of our statistical problem.
  \item $V$ is interpreted as potential energy, and is determined by the unextended target $f$.
  \item $K$ is interpreted as kinetic energy, and in HMC is part of the algorithmic design, a
    standard choice being $K(p) = \frac{1}{2}p^T M^{-1}p$ for some symmetric positive definite
    $M \in \R^{d\times d}$.
  \item $\mu^H$ is interpreted as a Gibbs measure that governs the distribution of $(q,p)$
    over an ensemble of many copies of the system in contact with a heat bath at a constant
    temperature.
\end{itemize}

Notice that $f$ is the marginal of $f^H$ over the position variables, and therefore we can
obtain samples from $f$ by sampling $f^H$ and throwing away the momentum variables. At first
sight, it is not clear how sampling $f^H$ rather than directly sampling $f$ can help. The main
idea is that extending the state space allows us to exploit several properties of Hamiltonian
dynamics, proposing bolder moves in state space while not hindering the acceptance rate and the
mixing of the chain.

%% Chapter 8 — Section 8.1: Hamiltonian Dynamics
%% Source: PDF pages 100–102 (book pages 87–89 + 90–91)

\section{Hamiltonian Dynamics}
\label{sec:8.1}

This chapter is organized as follows. Hamiltonian dynamics are reviewed in Section~\ref{sec:8.1},
emphasizing three properties that are important to the design of sampling algorithms. The HMC
algorithm and an idealized version of it are introduced in Section~\ref{sec:8.2}. We focus on
showing that these algorithms are exact, meaning that they leave the target distribution
invariant. Section~\ref{sec:8.3} contains bibliographical remarks, including references that
address the geometric ergodicity of HMC.

\medskip
Given a Hamiltonian $H(q,p)$, the corresponding Hamilton equations are
\begin{equation}
\label{eq:8.2}
\begin{aligned}
\frac{dq_i}{dt} &= \frac{\partial H}{\partial p_i}, & 1 \le i \le d, \\[4pt]
\frac{dp_i}{dt} &= -\frac{\partial H}{\partial q_i}, & 1 \le i \le d.
\end{aligned}
\end{equation}

We vectorize equations~\eqref{eq:8.2} as follows
\[
\frac{dx}{dt} = J^{-1}\nabla H(x), \qquad
J := \begin{bmatrix} 0 & -I_{d\times d} \\ I_{d\times d} & 0 \end{bmatrix}
\in \R^{2d\times 2d}, \quad
x := \begin{bmatrix} q \\ p \end{bmatrix} \in \R^{2d},
\]
where
\[
\nabla H = \left[\frac{\partial H}{\partial q_1}, \ldots,
\frac{\partial H}{\partial q_d},
\frac{\partial H}{\partial p_1}, \ldots,
\frac{\partial H}{\partial p_d}\right]^T, \qquad
\frac{d}{dt}\begin{bmatrix}q\\p\end{bmatrix}
= \left[\frac{dq_1}{dt},\ldots,\frac{dq_d}{dt},\frac{dp_1}{dt},\ldots,\frac{dp_d}{dt}\right]^T.
\]

The matrix $J$ plays a central role in the Hamiltonian formalism. Note that $J$ is
skew-symmetric, meaning that $J^T = -J$. Later we will use the following properties of
invertible skew-symmetric matrices:
\begin{enumerate}
  \item $J^{-1}$ is also skew-symmetric, since $J^{-1} = -(J^T)^{-1} = -(J^{-1})^T$.
  \item The corresponding quadratic form is zero, since $x^T J x = x^T J^T x = -x^T J x$.
\end{enumerate}

Hamilton equations~\eqref{eq:8.2} define a flow in phase space: For $t \ge 0$, we denote by
$\varphi_t : \R^{2d} \to \R^{2d}$ the $t$-flow map, so that $\varphi_t(x) \in \R^{2d}$ is
the solution to Hamilton equations at time $t$ with initial condition $x \in \R^{2d}$. It
follows directly from the definition that $\varphi_0(x) = x$, i.e.\ $\varphi_0$ is the
identity map. Furthermore, for $t, s \ge 0$ we have that $\varphi_{s+t}(x) = \varphi_s \circ
\varphi_t(x)$, since the solution at time $t + s$ with initial condition $x$ can be found by
first solving up to time $t$ to obtain $x(t) := \varphi_t(x)$, and then solving Hamilton
equations for an additional time $s$ with initial condition $x(t)$ to obtain
$x(t+s) := \varphi_s(x(t)) = \varphi_s \circ \varphi_t(x)$. The concept of flow map of a
dynamical system is helpful in exploring the qualitative behavior of numerical methods to solve
them, since one can think of numerical methods as approximations of the flow map. In particular,
we will see that in order to ensure that Hamiltonian Monte Carlo algorithms leave the extended
target invariant, it is essential to use numerical solvers that preserve key properties of the
flow map.

\begin{example}[Harmonic Oscillator]
Consider the case $d = 1$, $V(q) = \frac{q^2}{2}$, and $K(p) = \frac{p^2}{2}$. Then, the
Hamilton equations are
\[
\frac{dq}{dt} = p, \qquad \frac{dp}{dt} = -q.
\]
The general solutions have the form
\[
q(t) = r\cos(a + t), \qquad p(t) = -r\sin(a + t),
\]
for some constants $a$ and $r$, which shows that $\varphi_t$ is a clockwise rotation in the
$q-p$ plane.
\end{example}

Hamiltonian dynamics dovetail nicely with MCMC because they satisfy three important properties:
conservation of the Hamiltonian, symplecticity, and reversibility. We next explain these
properties and how they are relevant to the construction of MCMC algorithms.

%% -------------------------------------------------------
\subsection{Conservation of Hamiltonian}
\label{sec:8.1.1}
%% -------------------------------------------------------

By the chain rule and the fact that $J$ is skew-symmetric, we have that
\[
\frac{dH}{dt} = \nabla H^T \frac{dx}{dt} = \nabla H^T J^{-1} \nabla H = 0, \qquad
x = \begin{bmatrix}q\\p\end{bmatrix}.
\]

Recall the definition of $f^H$ in \eqref{eq:8.1}. The value of $f^H$ depends on $x = (q,p)$
only through $H$. Therefore, conservation of the Hamiltonian implies that if we start at any
$x_0 = (q_0, p_0)$ and move with Hamiltonian dynamics, we move along level sets of $f^H$. In
other words, we have $H(x_0) = H(\varphi_s(x_0))$ for all $s > 0$. This property, together
with the preservation of the canonical measure that we describe next, will allow us to use
Hamiltonian dynamics to propose bold moves that are still likely to be accepted.

\begin{figure}[htbp]
\centering
\includegraphics[width=0.60\textwidth]{figures/fig_08_1}
\caption{\textit{Movement along a level set of $f^H$ starting from $(q_0, p_0)$.}}
\label{fig:8.1}
\end{figure}

%% -------------------------------------------------------
\subsection{Symplecticity and Preservation of the Canonical Measure}
\label{sec:8.1.2}
%% -------------------------------------------------------

The flow map $x(s) \xrightarrow{\varphi_t} x(t+s)$ preserves volume. This is a direct
consequence of the symplecticity of Hamiltonian dynamics. We next recall the definition of
symplecticity and prove the symplecticity of Hamiltonian dynamics using a theorem by Poincar\'{e}.

\begin{definition}[Symplectic Maps]
A linear map $L : \R^{2d} \to \R^{2d}$ is symplectic if $L^T J L = J$. A continuously
differentiable map $\psi : \R^{2d} \to \R^{2d}$ is symplectic if its Jacobian is symplectic
everywhere. That is, if
\[
\psi'(x)^T J \psi'(x) = J, \qquad \forall x \in \R^{2d}.
\]
\end{definition}

\begin{remark}
In $\R^2$ symplecticity of linear maps is equivalent to area preservation of parallelograms;
in $\R^{2d}$ it is equivalent to conservation of the oriented areas of the projections of
parallelograms into the coordinate planes $(q_i, p_i)$.
\end{remark}

\begin{theorem}[Hamiltonian Flows are Symplectic]
Let $H$ be a $\mathcal{C}^2$ Hamiltonian. Then the corresponding flow map $\varphi_t$ is
symplectic for all $t$.
\end{theorem}

\begin{proof}
Let $\varphi(x,t) := \varphi_t(x)$ denote the flow map of time $t$ starting from $x$. We
need to show that, for any $t$ and arbitrary $x \in \R^{2d}$,
\begin{equation}
\label{eq:8.3}
\frac{\partial}{\partial x}\varphi(x,t)^T J \frac{\partial}{\partial x}\varphi(x,t) = J.
\end{equation}

First, the equality holds for $t = 0$, since $\varphi(x,0)$ is the identity map on $\R^{2d}$.
Therefore, the proof will be complete if we show that the left-hand side in
equation~\eqref{eq:8.3} is constant in time. To that end, recall that
$\frac{\partial}{\partial t}\varphi(x,t) = J^{-1}\nabla H(\varphi(x,t))$ and so, using that
$H$ is twice continuously differentiable, we have
\begin{align*}
\frac{\partial^2}{\partial t \partial x}\varphi(x,t)
&= \frac{\partial^2}{\partial x \partial t}\varphi(x,t) \\
&= \frac{\partial}{\partial x} J^{-1}\nabla H(\varphi(x,t)) \\
&= J^{-1}\nabla^2 H\!\left(\varphi(x,t)\right)\frac{\partial}{\partial x}\varphi(x,t),
\end{align*}
where $\nabla^2$ denotes the Hessian. This partial result and the skew-symmetry of $J$ give that
\begin{align*}
&\frac{\partial}{\partial t}\!\left(\frac{\partial}{\partial x}\varphi(x,t)^T J
  \frac{\partial}{\partial x}\varphi(x,t)\right) \\
&= \frac{\partial^2}{\partial t\partial x}\varphi(x,t)^T J
  \frac{\partial}{\partial x}\varphi(x,t)
  + \frac{\partial}{\partial x}\varphi(x,t)^T J
  \frac{\partial}{\partial t\partial x}\varphi(x,t) \\
&= \frac{\partial}{\partial x}\varphi(x,t)^T \nabla^2 H\!\left(\varphi(x,t)\right)
  J^{-T} J \frac{\partial}{\partial x}\varphi(x,t)
  + \frac{\partial}{\partial x}\varphi(x,t)^T J J J^{-1}\nabla^2 H\!\left(\varphi(x,t)\right)
  \frac{\partial}{\partial x}\varphi(x,t) \\
&= -\frac{\partial}{\partial x}\varphi(x,t)^T \nabla^2 H\!\left(\varphi(x,t)\right)
  \frac{\partial}{\partial x}\varphi(x,t)
  + \frac{\partial}{\partial x}\varphi(x,t)^T \nabla^2 H\!\left(\varphi(x,t)\right)
  \frac{\partial}{\partial x}\varphi(x,t) \\
&= 0,
\end{align*}
and the proof is complete.
\end{proof}

Taking the determinant at both sides of the equality
\[
\left(\varphi_t'\right)^T J \varphi_t = J
\]
shows that the Jacobian determinant of $\varphi_t$ has absolute value 1. Using this property
and the conservation of the Hamiltonian, we obtain the following important result:

\begin{theorem}[Preservation of Canonical Measure]
The flow $\varphi_t$ preserves the canonical measure of $H$.
\end{theorem}

\begin{proof}
Let $D$ be a Lebesgue measurable set. We have
\begin{align*}
\mu^H\!\left(\varphi_t^{-1}(D)\right)
&= \int_{\varphi_{-t}(D)} \frac{1}{Z} e^{-H(x)}\,dx \\
&= \int_D \frac{1}{Z} e^{-H(\varphi_{-t}(x))} \left|\varphi_{-t}'(x)\right| dx \\
&= \int_D \frac{1}{Z} e^{-H(x)}\,dx \\
&= \mu^H(D).
\qedhere
\end{align*}
\end{proof}

%% -------------------------------------------------------
\subsection{Time Reversibility}
\label{sec:8.1.3}
%% -------------------------------------------------------

The term time reversible comes from classical mechanics. Let $q(t)$ and $p(t)$ be the height
and momentum of a pendulum of mass $m$ at time $t$. Suppose we release the pendulum at
$q(0) = q_0$ with momentum $p(0) = v_0 m$ in a frictionless space. By conservation of energy,
we know that the pendulum will come back to the same position $q_0$ at some time $t$, and its
velocity will be $-v_0$. The $q-p$ time plot is symmetric i.e.\ the system is time reversible.
Here we give a definition of time reversibility in the context of dynamical systems.

\begin{definition}[Involution and Reversibility]
A linear map $S$ is an involution if $S^2 = S \circ S$ is the identity map. A bijection
$\varphi$ is called reversible with respect to an involution $S$ if $S \circ \varphi =
\varphi^{-1} \circ S$. A differential equation is called reversible with respect to an
involution $S$ if, for every fixed $t$, its flow $\varphi_t$ is reversible with respect to $S$.
\end{definition}

The inverse of the flow $\varphi_t$ of a differential equation $\frac{dx}{dt} = u(x)$ can be
intuitively viewed as the flow map $\varphi_t$ of the system $\frac{dx}{dt} = -u(x)$. We can
give an equivalent characterization of time reversibility:

\begin{theorem}[Characterization of Reversibility]
A system of differential equations $\frac{dx}{dt} = u(x)$ is time reversible with respect to
an involution $S$ if and only if $S \circ u = -u \circ S$.
\end{theorem}

\begin{proof}
Suppose $\frac{dx}{dt} = u(x)$ is time reversible with respect to $S$. Using linearity of
$S$, we have
\begin{align*}
S \circ u(x)
&= \lim_{\epsilon \to 0} \frac{1}{\epsilon}
  \Bigl(S \circ \varphi_{t+\epsilon}(x) - S \circ \varphi_t(x)\Bigr) \\
&= \lim_{\epsilon \to 0} \frac{1}{\epsilon}
  \Bigl(\varphi_{t+\epsilon}^{-1} \circ S(x) - \varphi_t^{-1} \circ S(x)\Bigr) \\
&= \frac{d}{dt}\varphi_t^{-1}\!\left(S(x)\right) \\
&= (-u) \circ S(x).
\end{align*}
The other direction follows directly from the remark above.
\end{proof}

\begin{corollary}[Reversibility of Hamiltonian Systems]
Hamiltonian systems with Hamiltonian $H(q,p) = V(q) + K(p)$ such that $K(p)$ is an even
function are reversible with respect to the momentum flip involution
$S : (q,p) \mapsto (q,-p)$.
\end{corollary}

\begin{example}[Gaussian Momentum Variables and Reversibility]
The kinetic energy $K(p) = \frac{1}{2}p^T M^{-1} p$ is an even function, since
\[
K(-p) = \frac{1}{2}(-p)^T M^{-1}(-p) = \frac{1}{2}p^T M^{-1} p = K(p).
\]
Therefore, Corollary~8.8 implies that Hamiltonian systems with Hamiltonian
$H(q,p) = V(q) + \frac{1}{2}p^T M^{-1}p$ are always time reversible with respect to
momentum flip. Notice that with the choice of kinetic energy $K(p) = \frac{1}{2}p^T M^{-1}p$,
the marginal of $f^H(q,p) \propto \exp(-H(q,p))$ over the momentum variables is
Gaussian $\mathcal{N}(0, M)$.
\end{example}

%% Chapter 8 — Section 8.2: From Hamiltonian Dynamics to a Sampling Algorithm
%% Source: PDF pages 105–107 (book pages 92–94)

\section{From Hamiltonian Dynamics to a Sampling Algorithm}
\label{sec:8.2}

%% -------------------------------------------------------
\subsection{Idealized Hamiltonian Monte Carlo Algorithm}
\label{sec:8.2.1}
%% -------------------------------------------------------

For the following algorithm, we consider a Hamiltonian of the form
$H(q,p) = \frac{1}{2}p^T M^{-1}p + V(q)$. We would like to construct a Markov kernel that
leaves the canonical measure of $H$ invariant. Recall that our true target is the $q$-marginal
with density proportional to $e^{-V}$; the Gaussian momentum $p$ is just an auxiliary variable.
As usual, we have access to the potential $V(q)$ (and its gradient) but not to the normalizing
constant $Z$.

\begin{algorithm}[H]
\caption{Idealized, Not-Implementable Hamiltonian Monte Carlo}
\label{alg:8.1}
\begin{algorithmic}[1]
\Require Initialization $(q^{(0)}, p^{(0)})$, duration parameter $\lambda$, sample size $N$.
\For{$n = 0, \ldots, N-1$}
  \State \textbf{Momentum refreshment:} Sample $p^* \sim \mathcal{N}(0, M)$.
  \State \textbf{State update:} $(q^{(n+1)}, p^{(n+1)}) = \varphi_\lambda(q^{(n)}, p^*)$.
\EndFor
\Ensure Sample $\{q^{(n)}\}_{n=1}^N$.
\end{algorithmic}
\end{algorithm}

Note that the idealized, not implementable HMC algorithm shown above presupposes that the flow
map $\varphi_\lambda$ can be evaluated, which is not true for most practical applications.
However, it is important to note that if the flow map could be evaluated, there would be no need
for an acceptance/rejection step to leave the target invariant; this is proved in
Theorem~\ref{thm:8.10}. In practice, $\varphi_\lambda$ is approximated with numerical
integrators (which might not preserve total energy) and the acceptance/rejection step is used to
counter that error. This will be made precise in Algorithm~\ref{alg:8.2} and
Theorem~\ref{thm:8.12}.

\begin{theorem}
\label{thm:8.10}
Let $\mu^H$ be the Gibbs distribution with energy $H(q,p) = \frac{1}{2}p^T M^{-1}p + V(q)$.
The Markov kernel implicitly defined by Algorithm~\ref{alg:8.1} leaves $\mu^H$ invariant.
\end{theorem}

\begin{proof}
By the definition of the Hamiltonian, the momentum variables are $\mathcal{N}(0,M)$ and
therefore $\mu^H$ is clearly invariant in the momentum refreshment step. On the other hand,
Theorem~8.5 shows that the flow map $\varphi_\lambda$ is measure-preserving with respect to
the canonical measure of $H$. This implies that step~2 also preserves the distribution $\mu^H$.
\end{proof}

%% -------------------------------------------------------
\subsection{Hamiltonian Monte Carlo Algorithm}
\label{sec:8.2.2}
%% -------------------------------------------------------

Time reversibility, conservation of Hamiltonian, and symplecticity are properties that hold for
Hamiltonian systems. In practice, for most problems of interest, Hamilton equations cannot be
solved in closed form and need to be discretized. Discretizations of Hamiltonian systems may
preserve \emph{exactly} some of these properties. For instance, the famous leapfrog method is
reversible and symplectic. We remark that unfortunately a numerical solver for Hamilton
equations cannot be simultaneously symplectic \emph{and} conserve the Hamiltonian exactly.

\begin{example}[The Leapfrog Integrator]
\label{ex:8.11}
The leapfrog method applied to Hamilton equations with
$K(p) := \sum_{i=1}^d \frac{p_i^2}{2m_i}$ and step-size $\epsilon > 0$ is
\begin{align*}
p_i\!\left(t + \tfrac{\varepsilon}{2}\right)
  &= p_i(t) - \frac{\varepsilon}{2}\frac{\partial V}{\partial q_i}(q(t))
  &\quad&\text{(half momentum step)} \\[4pt]
q_i(t + \varepsilon)
  &= q_i(t) + \varepsilon \frac{p_i\!\left(t + \frac{\varepsilon}{2}\right)}{m_i}
  &\quad&\text{(full position step)} \\[4pt]
p_i(t + \varepsilon)
  &= p_i\!\left(t + \tfrac{\varepsilon}{2}\right)
    - \frac{\varepsilon}{2}\frac{\partial V}{\partial q_i}(q(t+\varepsilon)).
  &\quad&\text{(half momentum step)}
\end{align*}
It is time reversible and conserves volume exactly. To discretize Hamilton equations for a time
interval $[t, t+\lambda]$ of length $\lambda$, we may take $L$ leapfrog steps with
$\lambda = \varepsilon L$.
\end{example}

\begin{figure}[htbp]
\centering
\includegraphics[width=0.95\textwidth]{figures/fig_08_2}
\caption{\textit{Simulate the Hamiltonian $\frac{3p^2}{2} + \frac{4q^2}{2}$ using leapfrog
and Euler method updates: $p(t+\varepsilon) = p(t) + \varepsilon\frac{dp}{dt}$,
$q(t+\varepsilon) = q(t) + \varepsilon\frac{dq}{dt}$.}}
\label{fig:8.2}
\end{figure}

Let $\Psi_\lambda$ be a symplectic numerical integrator, time reversible with respect to the
momentum flip involution, that approximates the flow map $\varphi_\lambda$. The Hamiltonian
Monte Carlo algorithm proposes moves using $\Psi_\lambda$ and corrects for the fact that
$\Psi_\lambda$ does not preserve the Hamiltonian by using an accept/reject mechanism.

\begin{algorithm}[H]
\caption{Hamiltonian Monte Carlo (HMC)}
\label{alg:8.2}
\begin{algorithmic}[1]
\Require Initialization $(q^{(0)}, p^{(0)})$, duration parameter $\lambda$, sample size $N$.
\For{$n = 0, \ldots, N-1$}
  \State \textbf{Momentum refreshment:} Sample $p^{\mathrm{new}} \sim \mathcal{N}(0, M)$.
  \State \textbf{Propose new state:} $(q^*, p^*) = \Psi_\lambda(q^{(n)}, p^{\mathrm{new}})$.
  \State \textbf{Set}
  \[
    (q^{(n+1)}, p^{(n+1)}) =
    \begin{cases}
      (q^*, p^*) & \text{with probability }
        a := \min\!\left\{1,\, e^{H(q^{(n)}, p^{\mathrm{new}}) - H(q^*, p^*)}\right\}, \\
      (q^{(n)}, -p^{\mathrm{new}}) & \text{with probability } 1 - a.
    \end{cases}
  \]
\EndFor
\Ensure Sample $\{q^{(n)}\}_{n=1}^N$.
\end{algorithmic}
\end{algorithm}

\begin{center}
\fbox{\parbox{0.92\textwidth}{%
\textbf{Pros and Cons}: Using Hamiltonian dynamics in an extended state space, HMC breaks away
from the local behavior of RWMH and MALA proposals. As a result, HMC scales better to
high-dimensional problems (see below for further discussion). However, HMC is often harder to
implement and to tune than RWMH and MALA. Similar to MALA, HMC requires evaluation of the
gradient of $V$, which can be expensive.}}
\end{center}

\medskip
\noindent\textbf{Scaling of HMC} \quad
Classical theoretical results for HMC developed in the same large-$d$ scaling setting that we
considered for RWMH in Chapter~4 and for MALA in Chapter~6 revealed the scaling of the
leapfrog step-size should be $\mathcal{O}(d^{-1/4})$. Thus, HMC requires $\mathcal{O}(d^{1/4})$
steps to traverse the state space, which should be compared to $\mathcal{O}(d^{1/3})$ for MALA
and $\mathcal{O}(d^1)$ for RWMH. These analyses also show that the optimal acceptance rate for
HMC is $0.651$, compared to $0.574$ for MALA and $0.234$ for RWMH.

Our goal in the rest of this section is to understand why the acceptance probability in the HMC
algorithm is the correct one to ensure that $f^H$ invariant is an invariant distribution. To
that end, we will make use of the time reversibility of the numerical integrator and its
preservation of volume (otherwise a Jacobian would appear in the acceptance probability). The
main theoretical result of this chapter is the following.

\begin{theorem}[HMC Leaves Target Invariant]
\label{thm:8.12}
Suppose that $\Psi_\lambda$ is symplectic and reversible with respect to the momentum flip
involution. Then, the Markov kernel implicitly defined by HMC leaves invariant the canonical
measure of $H$.
\end{theorem}

\begin{proof}
Throughout the proof, we denote by $S$ the momentum flip involution. First, notice that the
momentum refreshment step preserves the canonical measure $\mu^H$. Hence, it suffices to show
that $\mu^H$ is invariant under steps~2 and~3. The corresponding Markov kernel is given by
\[
\pi_x(dy) = a(x)\,\delta_{\Psi_\lambda(x)}(y)\,dy
+ \bigl(1 - a(x)\bigr)\,\delta_{S(x)}(y)\,dy,
\]
where
\[
a(x) := \min\!\left\{1,\, e^{H(x) - H(\Psi_\lambda(x))}\right\}.
\]
Let $A \subset \R^{2d}$ be a Lebesgue measurable set. We need to show that
\[
\int_{\R^{2d}} \pi_x(A)\,\mu^H(dx) = \int_A \mu^H(dx).
\]
Expanding the left-hand side gives
\begin{align*}
\int_{\R^{2d}} \pi_x(A)\,\mu^H(dx)
&= \int_{\R^{2d}} a(x)\,\mathbf{1}_A\!\left(\Psi_\lambda(x)\right) f^H(x)\,dx
  + \int_{\R^{2d}} \bigl(1-a(x)\bigr)\,\mathbf{1}_A\!\left(S(x)\right) f^H(x)\,dx \\
&= \int_{\R^{2d}} \mathbf{1}_A\!\left(S(x)\right) f^H(x)\,dx
  + \int_{\R^{2d}} a(x)\,\mathbf{1}_A\!\left(\Psi_\lambda(x)\right) f^H(x)\,dx
  - \int_{\R^{2d}} a(x)\,\mathbf{1}_A\!\left(S(x)\right) f^H(x)\,dx.
\end{align*}
First, $\int_{\R^{2d}} \mathbf{1}_A(S(x)) f^H(x)\,dx = \int_A \mu^H(dx)$ since $S$ preserves
the canonical measure. Next, we use a change of variables in the third term to show that it is
equal to the second one. Precisely, using that $S \circ S = \mathrm{id}$, and that the
determinant Jacobians of $\Psi_\lambda$ and $S$ are identically one due to the symplecticity of
$\Psi_\lambda$ and reversibility of $S$, we have
\[
\int_{\R^{2d}} a(x)\,\mathbf{1}_A\!\left(S(x)\right) f^H(x)\,dx
= \int_{\R^{2d}} a\!\left(S \circ \Psi_\lambda(x)\right)
  \mathbf{1}_A\!\left(\Psi_\lambda(x)\right) f^H\!\left(\Psi_\lambda(x)\right) dx.
\]
Now, using that $\Psi_\lambda \circ S \circ \Psi_\lambda = S$ since $\Psi_\lambda$ is time
reversible with respect to $S$ and also that $H \circ S = H$, we deduce that
\begin{align*}
a\!\left(S \circ \Psi_\lambda(x)\right)
&= \min\!\left\{1,\, e^{H(S \circ \Psi_\lambda(x)) - H(\Psi_\lambda \circ S \circ \Psi_\lambda(x))}\right\} \\
&= \min\!\left\{1,\, e^{H(\Psi_\lambda(x)) - H(S(x))}\right\} \\
&= \min\!\left\{1,\, e^{H(\Psi_\lambda(x)) - H(x)}\right\} \\
&= \frac{f^H(x)}{f^H\!\left(\Psi_\lambda(x)\right)}\cdot a(x).
\end{align*}
Substituting this back to the previous integral gives
\[
\int_{\R^{2d}} a(x)\,\mathbf{1}_A\!\left(S(x)\right) f^H(x)\,dx
= \int_{\R^{2d}} a(x)\,\mathbf{1}_A\!\left(\Psi_\lambda(x)\right) f^H(x)\,dx,
\]
as desired.
\end{proof}

%% Chapter 8 — Section 8.3: Discussion and Bibliography
%% Source: PDF page 108 (book page 95)

\section{Discussion and Bibliography}
\label{sec:8.3}

HMC was proposed in the physics literature under the name Hybrid Monte Carlo~\cite{60}. The
method was introduced to the statistics community through the work~\cite{171}. An accessible
introduction to HMC is~\cite{170}; see also~\cite{23}. Theorem~8.4, which shows that
Hamiltonian flows are symplectic, was proved in~\cite{187}. Further background on Hamiltonian
dynamics and their numerical solution can be found in~\cite{210,101}. Numerical integrators
tailored to HMC are reviewed and compared in~\cite{31,42}.

In this chapter, we have seen that the Metropolis Hastings framework can be generalized by
means of an involution map, which we take as a momentum flip for HMC methods. The use of
involutions within the Metropolis Hastings algorithm was proposed in~\cite{11,92}, and has been
leveraged to establish mixing rates for Hamiltonian Monte Carlo in~\cite{92}. The theory
in~\cite{92}, which relies on the weak Harris approach developed in~\cite{102}, also covers
implementations of HMC in infinite-dimensional Hilbert spaces~\cite{22}. The paper~\cite{30}
relies instead on a novel coupling technique to analyze the convergence of HMC algorithms.
Further convergence results can be found in~\cite{148,65} and optimal scaling of HMC in high
dimensions was studied in~\cite{21}. Tuning HMC algorithms can be challenging. The No-U-Turn
Sampler (NUTS) algorithm~\cite{113} introduces a computational framework to automatically set
the duration parameter $\lambda$ and the step-size $\epsilon$.

%% Chapter 8 — Section 8.4: Exercises
%% Source: PDF pages 109–110 (book pages 96–97)

\section{Exercises}
\label{sec:8.4}

\begin{enumerate}

\item[8.1] Prove Corollary~8.8.

\item[8.2] Implement HMC to sample a bivariate Gaussian target $\mathcal{N}(\mu, \Sigma)$.
Consider parameter values $\mu = (0,0)$, $\sigma_X = 1$, $\sigma_Y = 1$ and two different
values of the correlation $\rho = 0$ and $\rho = 0.8$. Provide plots showing:
\begin{enumerate}
  \item The value of the Hamiltonian along trajectories of the leapfrog integrator for various
    choices of the step-size $\epsilon$. The goal is to show that if $\epsilon$ is small these
    trajectories are close to level-sets from the Gaussian target, but the accuracy deteriorates
    if $\epsilon$ is too large.
  \item A scatterplot of samples from HMC with different choices of step-size $\epsilon$ and
    duration parameter $\lambda$.
\end{enumerate}

\item[8.3] Recall from Exercise~6.3 that a response random variable $Y$ given explanatory
variables $(X_1, \ldots, X_d)$ follows a logistics regression model if
\[
Y = \begin{cases} 1 & \text{with probability}\quad \alpha, \\ 0 & \text{with probability}\quad 1 - \alpha, \end{cases}
\]
where
\[
\alpha = \frac{1}{1 + \exp\!\left(-(w_0 + w_1 X_1 + \ldots + w_d X_d)\right)}.
\]
\begin{enumerate}
  \item \textbf{Setting and Preparation}: Same as in Exercise~6.3.
  \item \textbf{Hamiltonian Monte Carlo}: Adopt a Bayesian formulation, specifying a Gaussian
    prior $w \sim \mathcal{N}(0, \sigma^2 I)$.
  \begin{enumerate}
    \item[(b)] Based on the derivations in Exercise~6.3, write down the potential $V(w)$ needed
      to implement HMC.
    \item[(c)] Choose
      \[
        M = \begin{bmatrix} 10 & 0 & 0 \\ 0 & 1 & 0 \\ 0 & 0 & 1 \end{bmatrix}
      \]
      and kinetic energy $K(p) = p^T M p$. Write down the leapfrog integrator applied to the
      Hamiltonian $H(w,p) = V(w) + K(p)$. Plot leapfrog trajectories for different choices of
      step-size.
    \item[(d)] Implement HMC for posterior sampling, tuning the step-size and duration
      parameters. Plot the value of the computed posterior mean of $w$ against the number of
      HMC samples.
  \end{enumerate}
\end{enumerate}

\item[8.4] Consider a 100-dimensional multivariate Gaussian with independent coordinates, mean
vector $\mu = 0$, and variance vector $\sigma^2 = (0.01, 0.02, \ldots, 1)$.
\begin{enumerate}
  \item[(a)] Sample using RWMH and plot histograms for the coordinates with variance $0.98$,
    $0.99$, and $1$.
  \item[(b)] Sample using MALA tuning the step-size. Plot histograms for the coordinates with
    variance $0.98$, $0.99$, and $1$.
  \item[(c)] Sample using HMC with $\lambda = 150$ and tuning the step-size. Plot histograms
    for the coordinates with variance $0.98$, $0.99$, and $1$.
  \item[(d)] For each of the 100 coordinates, compute the sample mean and sample variance of
    the samples you obtained with RWMH, MALA, and HMC. For each method, draw scatter plots of
    the 100 sample means and variances.
\end{enumerate}

\end{enumerate}

